\subsection{THETA3-C15: transferring curve factors to ball covers, and the countability boundary}
\label{C15}

We work over $\overline{\Q}\subset\C$ unless otherwise indicated.  Write
$C\rightsquigarrow D$ when a connected finite \'{e}tale cover of the smooth
projective curve $C$ admits a nonconstant morphism to the smooth projective
curve $D$.  This relation is directed.  It is not permissible to reverse an
arrow merely because the two curves occur in a correspondence.

\paragraph{Curve-to-ball transfer.}
Let $Y$ be a smooth compact arithmetic ball surface with a surjective morphism
$f:Y\to C_0$.  If $C_0\rightsquigarrow D$, there is a connected finite
\'{e}tale cover $X\to Y$ and a surjective morphism $X\to D$.  Consequently
$\operatorname{Alb}(X)\to J(D)$ is surjective.  Every abelian quotient of
$J(D)$ is therefore a quotient of the Albanese of an arithmetic ball cover.

\emph{Proof.}
Choose a connected finite \'{e}tale cover $C'\to C_0$ and a nonconstant
map $C'\to D$.  Any connected component $X$ of $Y\times_{C_0}C'$ is finite
\'{e}tale over $Y$.  Its image in the connected variety $Y$ is nonempty,
open, and closed, so this projection is surjective.  The image of $X$ in
$C'$ is closed by properness.  It cannot be a point, since its image in
$C_0$ is all of $C_0$; thus it is $C'$.  Composition gives the asserted
surjection $h:X\to D$.  The universal covering of $X$ is again the ball,
and its lattice has finite index in that of $Y$, so arithmeticity persists.

The pullback of a nonzero holomorphic differential on $D$ is nonzero on
$X$: this can be checked at a point where $h$ is smooth.  Thus
$h^*:H^0(D,\Omega_D^1)\to H^0(X,\Omega_X^1)$ is injective.  The dual map
on the tangent spaces of Albanese varieties is surjective.  The resulting
homomorphism $\operatorname{Alb}(X)\to J(D)$ is therefore surjective.
This also proves the assertion about further quotients. \hfill$\square$

The relation $\rightsquigarrow$ is transitive.  Indeed, if
$C'\to C$ is finite \'{e}tale, $C'\to D$ is nonconstant, and
$D'\to D$ is finite \'{e}tale with $D'\to E$ nonconstant, take a connected
component of $C'\times_DD'$.  It is finite \'{e}tale over $C'$ and maps
surjectively to $D'$, hence to $E$.  This proves transitivity, including the
connectedness qualification, without identifying any of the maps as an
isomorphism.

\paragraph{The exact Bogomolov--Tschinkel direction.}
Let $C_0$ now be Schoen's curve $y^2=1-x^5$.  Schoen constructs the morphism
$Y\to C_0$ on pp.~133--134 of his 2014 paper.  Put
\[
 D_6:\quad w^6=z(1-z),
\]
with smooth projective normalization understood.  This curve has genus two.
Bogomolov--Tschinkel, \emph{Unramified correspondences}, Proposition~1.8,
p.~2, proves
\[
 C_0\rightsquigarrow D_6.
\]
The same statement is Proposition~1.5, p.~2, of their
\emph{On curve correspondences}.  Their statement applies to every
hyperelliptic curve of genus at least two over an algebraically closed field
of characteristic different from $2,3$.  It does not assert
$D_6\rightsquigarrow C_0$ or that $D_6$ is universal in characteristic zero.
Their universality theorem over $\overline{\mathbf F}_p$, $p\geq5$, is
Theorem~1.7 of the first paper and Theorem~1.4 of the second; that theorem
must not be transported to characteristic zero.

There is nevertheless a useful positive continuation.  In
\emph{Unramified correspondences}, Example~3.1, pp.~5--6, let
$E_*\to\mathbf P^1$ be the degree-three quotient of the elliptic curve
with $j=0$, branched over $0,1,\infty$.  The paper proves
$D_6\rightsquigarrow H$ for curves $H$ admitting a map to $\mathbf P^1$
with local ramification indices at most two and branch points contained in
the image of $E_*(\overline{\Q})_{\rm tors}$.  It also states density of
these families in the indicated Hurwitz spaces.  Transitivity and the
preceding transfer lemma give actual quotients $J(H)$ of Albaneses of
covers of the fixed arithmetic ball surface $Y$.  This is a composition
of arrows in their proved directions.  In C16 we give a direct version
starting from Schoen's non-CM elliptic curve, with a full ramification
proof, an explicit new factor, and infinitely many nonisogenous factors.

\paragraph{A precise conditional universality reduction.}
If $D_6$ is universal over $\overline{\Q}$ in the sense of
Bogomolov--Tschinkel, then every abelian variety over $\overline{\Q}$ is a
quotient of the Albanese of a connected finite \'{e}tale cover of $Y$.
To see this, let $A$ be an abelian variety.  There is a smooth projective
curve $H$ with a surjection $J(H)\to A$.  For $\dim A=1$ take $H=A$;
in larger dimension one may take a smooth sufficiently ample complete
intersection curve in $A$, whose map on $H^1$ is injective by weak
Lefschetz.  Universality gives $D_6\rightsquigarrow H$, and then
$C_0\rightsquigarrow H$.  Apply the transfer lemma and compose the two
surjective homomorphisms of abelian varieties.  The hypothesis here is a
characteristic-zero universality assertion, not a consequence of the
positive-characteristic theorem.  Even under that hypothesis, the conclusion
concerns $H^1$ source factors; it does not produce the higher-degree cycle
projection or spanning property required in C04 and C13.

\paragraph{Countability.}
For a fixed compact ball surface $Y$, only countably many simple complex
abelian isogeny classes occur in the Albaneses of its finite \'{e}tale
covers.  In particular, for each $g\geq1$, a very general simple complex
abelian variety of dimension $g$ is not such a factor.

\emph{Proof.}
The compact manifold $Y$ has finitely generated fundamental group.  For
each $n$, it has only finitely many subgroups of index $n$: each such
subgroup gives a transitive permutation action on $n$ letters, and the
images of a fixed finite generating set admit only finitely many choices
in the finite group $S_n$.  Thus the connected finite covers form a
countable set.  Each Albanese has finitely many simple isogeny factors,
by complete reducibility of abelian varieties.  A countable union of
finite sets remains countable.

For a fixed complex abelian variety, there are countably many lattices
commensurable with its period lattice.  Thus its isogeny class is
countable, and allowing polarizations still gives a countable set because
integral alternating forms on a fixed lattice are countable.  The
positive-dimensional moduli space of polarized abelian varieties of
dimension $g$ has uncountably many complex points; the countable set just
obtained is contained in a countable union of proper closed points.
The same exclusion holds on its simple locus.  Finally, a nonzero rational
Hodge map from the $H^1$ of a simple abelian variety to $H^1(X)$ would,
by the equivalence for polarizable Hodge structures of weight one, make
that variety an isogeny factor of $\operatorname{Alb}(X)$.  This gives the
claimed source obstruction. \hfill$\square$

The same proof applies to any countable collection of arithmetic towers.
In particular, the number-field Hermitian data and finite adelic levels
used in the present unitary constructions give such a collection: number
fields, finite matrices over them, and finite-level choices are countable.
This countability obstruction does not conflict with the conditional
statement over $\overline{\Q}$, where the set of algebraic varieties is
already countable.  Nor does countability preclude density; the explicit
families in the next subsection are countable and dense in the relevant
complex hyperelliptic moduli spaces.

\paragraph{Framework anchor.}
The transfer lemma uses actual curve maps and their Albanese realizations,
the weight-one case of the correspondence framework in Grothendieck's
letter of 12 August 1964, cote~44, scans~7--8.  It supplies source factors;
the existence of higher-degree algebraic projectors remains the separate
issue visible in cote~16, scans~44--46.

\subsection{THETA3-C16: infinitely many non-CM elliptic factors from torsion branch points}
\label{C16}

Schoen's construction supplies a connected curve $C$ with
\[
 C\longrightarrow C_0 \quad\hbox{finite \'{e}tale of degree }2,
 \qquad C\longrightarrow F \quad\hbox{of degree }2,
\]
where $C_0:y^2=1-x^5$ and
$F:s^2=x^4+x^3+x^2+x+1$; see his pp.~133--134.  The second double cover
is ramified.  Indeed, in the function field of $F$ it is obtained by
adjoining $\sqrt{1-x}$.  The two points over $x=1$ and the two points at
infinity have odd valuation of $1-x$, so they are branch points.
Choose one branch point $q_0$ as origin on $F$.  This does not change its
$j$-invariant, which Schoen computes as $2048$.

\paragraph{The torsion-branch transfer lemma.}
Let $E$ be any elliptic curve in characteristic zero, let $C\to E$ be a
connected double cover ramified at its origin, and put
$\pi:E\to E/\{\pm1\}\simeq\mathbf P^1$.  For an odd positive integer
$N$, define
\[
 S_N=\pi(E[N]),\qquad |S_N|=(N^2+1)/2.
\]
Suppose that $H$ is a smooth connected projective curve and
$\sigma:H\to\mathbf P^1$ is a finite morphism of degree $d$, all of whose
local ramification indices are at most two, with branch locus contained
in $S_N$.  Then a connected finite \'{e}tale cover $D\to C$ of degree
at most $dN^2$ admits a surjection $D\to H$ of degree at most four.

\emph{Proof.}
Form $C_N=C\times_{E,[N]}E$.  Its projection to $C$ is finite \'{e}tale
of degree $N^2$.  It is a smooth connected curve: as a double cover of
the second copy of $E$ it is still branched, at every point of $E[N]$,
and a split degree-two cover would have no branch points.  Its composite
map $a:C_N\to E\xrightarrow{\pi}\mathbf P^1$ has degree four.
Every local ramification index of $a$ over $S_N$ is even.  Namely, above
$P\in E[N]$ the map $C_N\to E$ has ramification index two; multiplication
by $N$ is \'{e}tale and carries $P$ to the specified branch point.
The further index of $\pi$ is either one or two.  Oddness of $N$ also
shows that the only fixed point of $[-1]$ in $E[N]$ is the origin,
which proves the displayed cardinality of $S_N$.

Normalize $C_N\times_{\mathbf P^1}H$ and choose a connected component
$D$.  Its projection to $C_N$ is \'{e}tale.  Here is the complete local
check.  Where $\sigma$ is unramified, base change is unramified.  At a
ramified point of $\sigma$, choose formal parameters so that its map is
$z=u^2$, after absorbing a unit into $u$.  The other map has the form
$z=v^{2e}$ after absorbing a unit into $v$, since its ramification index
is even and the residue field is algebraically closed of characteristic
zero.  The normalization of $u^2=v^{2e}$ is the disjoint union of the
two branches $u=\pm v^e$, each unramified over the $v$-disc.  This proves
the assertion also at points where several unramified sheets of $\sigma$
lie over the same branch value.  Therefore every connected component is
finite \'{e}tale over $C_N$, with degree at most $d$.  It maps onto $H$:
both maps in the original fiber product are finite and surjective, and
each component of its normalization has dimension one.  Its degree over
$H$ is at most the total base-change degree four.  Composing with
$C_N\to C$ proves both degree bounds. \hfill$\square$

Applied to Schoen's $C$, C15 gives a connected arithmetic ball cover
$X\to Y$ of degree at most $2dN^2$ whose Albanese surjects onto $J(H)$.
For any set $B\subset S_N$ of cardinality $2g+2$, its hyperelliptic
double cover $H_B\to\mathbf P^1$ satisfies the hypotheses, and hence
\[
 \deg(X/Y)\leq4N^2,\qquad
 \operatorname{Alb}(X)\twoheadrightarrow J(H_B).
\]
As $N$ ranges through odd integers, these $H_B$ are dense in the complex
hyperelliptic moduli space of each fixed genus $g\geq1$.  Indeed, odd-order
torsion is dense in the complex torus $E(\C)$, so its image in
$\mathbf P^1(\C)$ is dense.  Any finite ordered configuration of distinct
points can be approximated by distinct points in this image.  The finitely
many odd orders occurring in such an approximation divide a common odd
$N$.  Passing continuously from configurations to the moduli of their
double covers proves density.  The curves are all defined over
$\overline{\Q}$ when $E$ is.  Density by itself does not prove that a
specified countable family avoids the CM locus.  The following arguments
give that additional conclusion explicitly.

\paragraph{An explicit new non-CM isogeny class, with cover degree at most $36$.}
The elliptic curve
\[
 E:\quad y^2=x^3-30x-25
\]
has $j(E)=2048$ and is therefore isomorphic over $\overline{\Q}$ to
Schoen's $F$ with the chosen origin.  Its nonzero three-torsion points have
four distinct $x$-coordinates, the roots of
\[
 \phi(r)=r^4-60r^2-100r-300.
\]
For completeness, the usual duplication formula yields this equation
without requiring a division-polynomial table: the condition
$x(2P)=x(P)$, with $y(P)\neq0$, gives
$3x^4+6Ax^2+12Bx-A^2=0$ on $y^2=x^3+Ax+B$; it is exactly the
nonzero three-torsion condition because $2P=P$ is impossible for such a
point.  Put $A=-30$, $B=-25$ and divide by three.

Fix a root $r$ of $\phi$, omit it from the four finite branch points, and
include infinity.  The resulting elliptic curve $H_r$ is
\[
 y^2=x^3+rx^2+(r^2-60)x+(r^3-60r-100).
\]
Its cubic is $\phi(x)/(x-r)$, so it has three distinct roots.  After
shifting $x$ by $-r/3$, its short Weierstrass coefficients are
\[
 A_r=\frac23r^2-60,\qquad
 B_r=\frac{20}{27}r^3-40r-100.
\]
Substitution in $j=1728\cdot4A_r^3/(4A_r^3+27B_r^2)$ and reduction
modulo $\phi(r)$ give
\[
 j(H_r)=-\frac{512}{405}
       (r^3+24r^2-130r-1200).
\]
Each of the four values is a root of the primitive polynomial
\begin{align*}
 P(J)={}&19683J^4-40310784J^3-5096079360J^2\\
        &+37203006717952J+6412351813189632.
\end{align*}
Both identities are polynomial identities over $\Q$ after reduction by
$\phi$; the accompanying standard-library certificate
\texttt{check\_C16\_torsion\_family.py} verifies them using exact rational
arithmetic, rather than floating-point root approximations.

Here is an explicit irreducibility certificate.  The monic reduction of
$P$ modulo $11$ is
\[
 P_0(J)=J^4-2J^3-5J^2+5J-5.
\]
In $\mathbf F_{11}[J]/(P_0)$ one has
\[
 J^{121}=5J^3-3J^2-5J+4,
 \qquad J^{14641}=J,
\]
and in $\mathbf F_{11}[J]$ one has the exact identity
\begin{align*}
 1={}&(-5J^2-5J-3)P_0(J)\\
 &+(J^3+4J^2-5J+2)(5J^3-3J^2-6J+4).
\end{align*}
The identity $J^{14641}=J$ says that all irreducible factors of $P_0$
have degree dividing four; the Bezout identity excludes factors of degree
one or two.  Thus $P_0$ is irreducible, and so is $P$ over $\Q$.
Since $P$ is primitive with leading coefficient $19683=3^9\neq1$,
its roots are not algebraic integers.  Every $H_r$ is therefore non-CM.
Moreover, it is not isogenous to $E$.  Indeed, a cyclic isogeny gives a
monic integral modular equation in the target $j$-invariant, so integrality
of the source $j$ implies integrality of the target; any isogeny factors
into cyclic isogenies.  Since $j(E)=2048$ is integral, the assertion follows.
The same argument applies place by place.  The standard source statements
are Milne, \emph{Modular Functions and Modular Forms}, version~1.31,
Theorem~6.1 and its proof, pp.~88--90, the cyclic-isogeny interpretation
in Lemma~8.5 and the following discussion, p.~99, and CM integrality in
Corollaries~12.15--12.16, p.~125.

All four branch points of $H_r$ belong to $S_3$, so the torsion-branch
lemma and C15 give the claimed ball cover of degree at most
$4\cdot3^2=36$.  This provides a non-CM elliptic isogeny class different
from Schoen's original one inside the same arithmetic covering tower.
No assertion that the four $H_r$ are mutually nonisogenous is needed.

\paragraph{Infinitely many pairwise nonisogenous non-CM factors.}
There are infinitely many pairwise nonisogenous non-CM elliptic curves
$E_1,E_2,\ldots$ over $\overline{\Q}$, each a quotient of the Albanese of
a connected finite \'{e}tale cover of Schoen's fixed ball surface $Y$.
For every finite initial set, one common connected cover of $Y$ has
$E_1\times\cdots\times E_m$ as an Albanese quotient.

\emph{Proof.}
We first produce a usable ordinary prime outside any prescribed finite
set.  The discriminant of $E:y^2=x^3-30x-25$ is
$2^4 3^6 5^3$.  If $p\equiv1\pmod{20}$, then $p\geq41$, the reduction
is good, and all its two-torsion is rational over $\mathbf F_p$:
\[
 x^3-30x-25=(x+5)(x^2-5x-5),
\]
whose remaining roots are $(5\pm3\sqrt5)/2$.  The field $\mathbf F_p$
contains the fifth roots of unity, and hence $\sqrt5$, because
$p\equiv1\pmod5$.  Consequently $4$ divides $\#E(\mathbf F_p)$ and
\[
 a_p=p+1-\#E(\mathbf F_p)\equiv2\pmod4.
\]
The Hasse bound gives $|a_p|\leq2\sqrt p<p$, so $a_p\not\equiv0\pmod p$.
This implies ordinary reduction.  One can check the last implication
directly: if $\varphi$ is Frobenius, its dual satisfies
$\widehat\varphi=[a_p]-\varphi$ and has degree $p$; its differential is
the nonzero scalar $a_p$.  It is therefore separable of degree $p$.
Since $[p]=\widehat\varphi\varphi$ and Frobenius is bijective on geometric
points, $\overline E[p](\overline{\mathbf F}_p)$ has $p$ points.
The facts used here are explicitly stated and proved in Stoll,
\emph{Elliptic Curves}, summer~2023 notes, Lemma~11.16, p.~42,
Theorem~12.1, p.~44, and Theorem~13.2, pp.~48--49.

There are infinitely many primes $p\equiv1\pmod{20}$, with an elementary
proof sufficient here.  Given a finite set of primes, choose a positive
integer $M$ divisible by $20$ and by every prime in that set.  Any prime
$p$ dividing
\[
 \Phi_{20}(M)=M^8-M^6+M^4-M^2+1>1
\]
is outside the set and is different from $2,5$.  The identity
$(M^2+1)\Phi_{20}(M)=M^{10}+1$ implies $M^{10}=-1$ modulo $p$.
The order of $M$ modulo $p$ is therefore four or twenty.  In the former
case $M^2=-1$, which gives $\Phi_{20}(M)=5$ modulo $p$, a contradiction.
Thus the order is twenty and $p\equiv1\pmod{20}$.

Fix such a prime and an embedding $\overline{\Q}\hookrightarrow
\overline{\Q}_p$.  The reduction map on $p$-torsion has a kernel of
cardinality $p$ and image $\overline E[p](\overline{\mathbf F}_p)$ of
cardinality $p$.  We spell out why no lifting assumption is hidden here.
The finite flat group scheme $E[p]$ over a good-reduction discrete
valuation ring has rank $p^2$.  After a finite unramified extension its
$p$ distinct special-fiber points are rational.  Their idempotent
decomposition lifts over this henselian ring.  Each component has rank
$p$, since translations make all special-fiber local lengths equal.
Its characteristic-zero generic fiber is \'{e}tale with $p$ geometric
points.  Every special point therefore lifts, and exactly $p$ generic
points reduce to the origin.

Choose a nonzero $P\in E[p]$ reducing to the origin, and choose
$Q,R\in E[p]$ whose reductions are nonzero and distinct modulo sign.
There are $(p-1)/2\geq2$ available nonzero sign classes.  The coordinates
satisfy
\[
 v_p(x(P))<0,\qquad
 v_p(x(Q)),v_p(x(R))\geq0,\qquad
 v_p(x(R)-x(Q))=0.
\]
For the first inequality, the integral local parameter $t=-x/y$ at the
origin gives $x=t^{-2}$ times an integral unit, and $v_p(t(P))>0$.
The others follow by reduction at finite distinct $x$-coordinates.

Let $E'_p$ be the double cover of $\mathbf P^1$ branched at
$\infty,x(P),x(Q),x(R)$.  Its branch cross ratio can be chosen as
\[
 \lambda=\frac{x(R)-x(Q)}{x(P)-x(Q)},\qquad v_p(\lambda)>0.
\]
The Legendre formula gives
\[
 j(E'_p)=256\frac{(1-\lambda+\lambda^2)^3}
                         {\lambda^2(1-\lambda)^2},
 \qquad v_p(j(E'_p))=-2v_p(\lambda)<0.
\]
This is a non-CM elliptic curve, by CM integrality.  All its branch points
lie in $S_p$, so the torsion-branch lemma gives a ball cover of degree at
most $4p^2$ whose Albanese surjects onto it.

Construct the sequence inductively.  After finitely many $E_i$ have been
chosen, the set of rational primes at which any conjugate of one of their
$j$-invariants is nonintegral is finite: each $j(E_i)$ is an algebraic
number.  Choose the next prime $p\equiv1\pmod{20}$ outside that set and
apply the construction.  All earlier $j(E_i)$ are integral at every
place above this new $p$, whereas $j(E'_p)$ is nonintegral at the chosen
place.  The local form of the monic modular-equation argument therefore
excludes an isogeny to every earlier $E_i$.  It also excludes an isogeny
to Schoen's original integral-$j$ elliptic curve.  This proves the infinite
assertion.

Finally, any finite list of the resulting connected finite \'{e}tale
ball covers has a common connected finite \'{e}tale cover, obtained by
taking a connected component of their fiber product over $Y$.  Projection
to each connected factor is surjective.  The maps to the $E_i$ persist,
and their distinct simple isogeny classes imply that the direct sum of
their $H^1$ injects into the $H^1$ of this common cover.  By complete
reducibility, the product $E_1\times\cdots\times E_m$ is an Albanese
quotient. \hfill$\square$

\paragraph{Nonzero periods with two of the new factors.}
The two independent curve projections of C12 give a generically finite
map $Y\to C_0^2$ of degree $125$.  For two curves $H_1,H_2$ furnished by
the torsion-branch lemma, choose the corresponding connected covers
$D_i\to C_0$; each $D_i\to H_i$ has degree at most four.  A connected
component $X$ of
\[
 Y\times_{C_0^2}(D_1\times D_2)
\]
is a finite \'{e}tale arithmetic ball cover of $Y$ and maps dominantly to
$H_1\times H_2$, with degree
\[
 1\leq d\leq125\cdot4\cdot4=2000.
\]
Indeed, its degree over $D_1\times D_2$ is at most the total generic
base-change degree $125$, and the remaining product map has degree at
most $16$.  If the $H_i$ are elliptic, choose differentials $\alpha_i$
normalized by $i\int_{H_i}\alpha_i\wedge\overline\alpha_i=1$.
Writing $a_i:X\to H_i$ for the two maps, the projection formula yields
\[
 \int_X a_1^*\alpha_1\wedge a_2^*\alpha_2
       \wedge\overline{a_1^*\alpha_1}
       \wedge\overline{a_2^*\alpha_2}=d>0.
\]
Thus one can obtain fourfold periods with factors from infinitely many
non-CM isogeny classes, with a positive integral degree in each instance.
The map to $H_1\times H_2\times H_1\times H_2$ repeats the two
projections; its image cycle is $d$ times the product of the two diagonals.
These are divisor-generated classes.  The construction substantially
enlarges the non-CM $H^1$ sources, but C14 still prevents these particular
product-of-curves image cycles from detecting the generic exceptional
Weil component.  The geometric factors constructed here are not supplied
by the CM character packets in the congruence $H^1$ statement of C01.

\paragraph{Framework anchor.}
The new Albanese quotients and image surfaces are actual algebraic
correspondences, as required by the realization/polarization viewpoint
of cote~44, scans~7--8.  Their explicit divisor-generated image classes
must be distinguished from the additional algebraic inverse and projector
classes discussed in cote~16, scans~43--46.  The infinite family enlarges
the available non-CM realizations without resolving that higher-degree
existence question.
